\documentclass[11pt]{article}
\usepackage[margin=0.85in]{geometry}
\usepackage[T1]{fontenc}\usepackage{lmodern,microtype}
\usepackage{amsmath,amssymb,booktabs,longtable,array,graphicx,xcolor}
\usepackage[hidelinks]{hyperref}\usepackage{fancyhdr}
\definecolor{navy}{HTML}{173957}
\pagestyle{fancy}\fancyhf{}\lhead{Fixed-area gamma Moran study}\rhead{Four-year joint inference}\cfoot{\thepage}
\setlength{\headheight}{14pt}\setlength{\parindent}{0pt}\setlength{\parskip}{6pt}\setlength{\emergencystretch}{2em}
\title{\color{navy}Four-year clonal evolution after CAR-T:\\fixed-area gamma fitness, joint inference\\and stochastic Moran simulations}
\author{Confirmed study protocol and reproducible analysis}\date{5 November 2026}
\begin{document}\maketitle
\textbf{Confirmed correction.} The additional mutant fitness has a \emph{fixed area}, not a fixed peak. With time in days, the reference area is 1 day, implementing the original expression $1+\mathrm{GammaPDF}(t)$ with explicit units. Shape and scale change the timing, width and height while preserving the integrated additional advantage. Peak-normalized results from the earlier study are superseded for this analysis.

\textbf{Objectives.} Fit gamma shape/scale jointly to longitudinal VAF and source-supported tMN outcomes; repeat under fixed-parameter sensitivity scenarios; estimate year-1, -2, -3 and -4 cumulative model risks; and simulate 200 full stochastic Moran trajectories for each prescribed shape/scale and initial-double-hit setting. The requested execution order is code, TeX protocol, then computation.

\section{Three-type Moran population}
Let $W$, $C$ and $D$ be mutually exclusive wild-type, single-mutant and double-hit progenitor populations, with
\[
W+C+D=N=100000.
\]
$D$ aggregates descendants acquiring a second transforming alteration and retains the founder marker. Multiple independent second-hit origins can be aggregated because their fitness is identical. The effective competing population and two-hit mechanism are modeling assumptions, not universal molecular definitions of tMN.

CAR infusion is day 0. Before day 0,
\[
f_W=f_C=f_D=1.
\]
This removes LD-associated selection and pre-CAR double-hit advantage. Neutral turnover and mutation can still occur. The study does not introduce a separate LD population bottleneck.

Define a shape--scale gamma density
\[
g(t;\alpha,\theta)=\frac{t^{\alpha-1}e^{-t/\theta}}{\Gamma(\alpha)\theta^\alpha},
\quad t>0,\quad \alpha>1,\quad\theta>0.
\]
Here $\alpha$ is dimensionless, $\theta$ is a scale in days, and $g$ has units day$^{-1}$. After infusion,
\[
f_W=1,\qquad f_C(t)=1+A_{\mathrm{area}}g(t;\alpha,\theta),\qquad f_D(t)=f_C(t)+s_D.
\]
Set $g(0)=0$ for $\alpha>1$. Thus $D$ inherits the same gamma pulse and adds a fixed extra advantage after day 0. The peak day is $(\alpha-1)\theta$, its height is derived, and
\[
\int_0^\infty[f_C(t)-1]dt=A_{\mathrm{area}}=1\ \mathrm{day}.
\]
No peak-height normalization is applied anywhere in the confirmed study.
The area inside the observed four years is $A_{\mathrm{area}}\Pr\{G\le1460\}$ for $G\sim\mathrm{Gamma}(\alpha,\theta)$ and may be smaller than the full area if a pulse is placed beyond follow-up. Record this quantity separately.

Replacement opportunities occur at total rate $Nb$. If $F=W+f_CC+f_DD$, newborn probabilities are
\[
p_W=\frac WF,\qquad p_C=\frac{(1-\mu)f_CC}{F},\qquad p_D=\frac{f_DD+\mu f_CC}{F}.
\]
The off-diagonal generator is $q_{n,n+e_j-e_k}=bn_kp_j$ for $j\ne k$. Mutation is a probability per $C$ division, not per day. Wild-type-to-founder mutation is not included.

\section{Reference parameters and sensitivity}
\begin{center}\begin{tabular}{@{}lp{0.49\textwidth}l@{}}\toprule
Parameter & Definition & Reference\\\midrule
$N$ & Effective competing population & $10^5$\\
$b$ & Moran turnover scale, day$^{-1}$ & $1/28$\\
$s_D$ & Extra post-CAR double-hit fitness & 0.02\\
$A_{\mathrm{area}}$ & Integrated extra founder advantage, days & 1\\
$\mu$ & Second transforming event per $C$ division & $10^{-6}$\\
$D_{\mathrm{init}}$ & Clinical-fit double-hit count at baseline & 0\\
$\kappa$ & Strict aggregate event threshold fraction & 0.01\\
$\sigma_{\mathrm{VAF}}$ & Fixed logit observation SD & 0.25\\
Delay & Threshold-to-diagnosis delay & 0 days\\\bottomrule
\end{tabular}\end{center}
The reference $b$ has precedent in a published hematopoietic Moran framework, which also discusses slower self-renewal rates; it is not a universally measured progenitor division rate. See \href{https://pubmed.ncbi.nlm.nih.gov/33636128/}{Heyde et al., Cell (2021)}. The advantage and transforming-mutation values remain sensitivity assumptions. The previous rate pair $b=1/7$, $s_D=0.021$ is retained as a rate comparator under the new fixed-area convention.

Only one cohort-wide $(\alpha,\theta)$ pair is fitted in each scenario. Initially search $1.01\le\alpha\le50$, $0.25\le\theta\le1460$ days and report boundary optima. Fixed-parameter sensitivities are
\begin{align*}
b&\in\{1/280,1/28,1/7\}\ \mathrm{day}^{-1}, &s_D&\in\{0,0.01,0.02,0.05\},\\
A_{\mathrm{area}}&\in\{0.3,1,3\}\ \mathrm{days}, &\mu&\in\{0,10^{-7},10^{-6},10^{-5},10^{-4}\},\\
D_{\mathrm{init}}&\in\{0,1,10,100\}, &\kappa&\in\{0.005,0.01,0.02\},\\
\sigma_{\mathrm{VAF}}&\in\{0.1,0.25,0.5\}.&&
\end{align*}
Changing $b$ holds $s_D$ fixed; derived $r_D=bs_D$ changes. Area is fixed across alpha/theta values within each scenario; area sensitivity varies that fixed assumption between scenarios. Set $K=\lfloor\kappa N\rfloor+1$. Compatible scenarios refit the complete joint likelihood. With $\mu=D_{\mathrm{init}}=0$, observed tMN is structurally impossible and no probability floor is used.

\section{Clinical cohort and observation model}
Use the supplied MOESM7/MOESM9 source workbooks and the audited exact-locus longitudinal extraction. Preserve hashes, source rows, sample identities and allele orientation. The confirmed analysis has 14 patients, 42 follow-up VAF observations and one baseline-selected index marker per patient. Sample-specific PBL/PBMC data are kept separate; cfDNA and tumor are not pooled into the blood-cell observation model. Other variants are retained as source evidence but are not treated as independent patient outcomes.

The index marker is selected from a positive baseline at or before infusion using baseline information only. Condition on the measured baseline to initialize total founder burden, $C+D=2Nv_0$, and subtract $D_{\mathrm{init}}$ to initialize $C$. Baseline measurement error is not propagated in this initial observation specification. Assume a representative diploid heterozygous marker with
\[
E[VAF(t)\mid X(t)]=\frac{C(t)+D(t)}{2N},\qquad
\operatorname{logit}(2VAF_{\rm obs})\sim N\!\left(\operatorname{logit}\frac{C+D}{N},\sigma_{\mathrm{VAF}}^2\right).
\]
Positive VAF must be below 50\%. Non-detection needs a detection model; absent filtered calls are not zero and no read counts are invented. The positive-detection restriction, near-50\% transformation, potential LOH/CNV, and unphased founder choice remain limitations. Exclude observations at or after tMN diagnosis from the prediagnosis process.

\begin{center}\begin{tabular}{llp{0.47\textwidth}}\toprule
Patient & Outcome used & Source-record assumption\\\midrule
103 & Exact event, day 127 & Listed time is first tMN from the relevant index CAR infusion.\\
197 & Exact event, day 589 & Same assumption; include its prediagnosis VAF.\\
25 & Right censor, day 1460 & Listed first event at day 1561 establishes event-free status through four years.\\
Other 11 & Unknown outcome & Contribute VAF only; not interpreted as clinical negatives.\\\bottomrule
\end{tabular}\end{center}
Baseline-only patients with informative outcomes are included. Do not require all patients to have complete outcome verification before a joint fit. Patients 18/69 have unresolved time origins, 127 lacks an eligible pre-CAR blood-cell initial state, and 275 lacks diagnosis timing; retain these issues explicitly outside the primary cohort.

These source-record assumptions were authorized for exploratory joint analysis. Outcome-dependent ascertainment or missingness may bias an unadjusted fit; no claim of population-incidence calibration follows from using the available events. Mortality is not modeled as a competing event. The index marker need not be the actual clinical tMN founder, which is another substantive approximation.

\section{Shared-state joint likelihood and numerical method}
With $Y_i$ the longitudinal VAF and $O_i$ the clinical observation,
\[
L(\alpha,\theta\mid\psi)=\prod_i\int p(X_i\mid\alpha,\theta,\psi)
\,p(Y_i\mid X_i,\psi)\,p(O_i\mid X_i,\psi)\,dX_i.
\]
Equivalently, $p(Y_i)\,p(O_i\mid Y_i)$; do not multiply independent VAF and clinical marginals. Apply VAF likelihood weights sequentially to the same latent-state distribution, and use boundary flux for exact tMN time or survival for censoring. Unknown outcomes have no clinical observation factor. Each patient's outcome is used once, without arbitrary VAF/event reweighting.

For tractable inference, retain a deterministic background $q=C/(N-D)$,
\[
\dot q=bq(1-q)\frac{(1-\mu)f_C-1}{1+(f_C-1)q},
\]
and a stochastic $D$ chain with conditional $C=(N-D)q$, birth $b(N-d)p_D$ and death $bd(1-p_D)$. This approximates stochastic single-mutant variation and feedback. Full W/C/D stochastic simulations below assess its behavior; they are not replaced by this approximation.

Known events/censoring use a killed chain below $K$. Unknown outcomes use an unrestricted physical chain with a numerical reflecting cap; compare cap refinements and, for the reference, the full $D=0,\ldots,N$ domain. Reflection at a numerical cap must not be interpreted as no tMN. Control changes in log likelihood, not only lost probability mass.

To resolve tiny early-event probabilities, evolve $v_d=e^{\eta(t)d}p_d$ under the similarity-transformed generator, including the additional $\dot\eta(t)d$ diagonal term, with explicit absorbed flux and rescaling. A branching-process tilt schedule stabilizes the growing rare-event tail; this is a numerical change of variables, not a fitted biological parameter or a change to the finite-N rates. Transform back in log space for event density, survival and VAF weighting. Check against an independently exponentiated generator; do not replace small event probabilities by an epsilon. Adaptive banded LSODA and a separately solved background are tolerance-refined.
When a narrow VAF error model can amplify extremely small D-tail errors, apply the same invertible tilt to the unrestricted unknown-outcome chain. Refine the state cap and solver tolerances before accepting those likelihoods. Limit the initial numerical tilt for large starting D populations to preserve physical mass.

Optimize from multiple starting points. Report joint NLL, marginal VAF NLL and clinical-given-VAF NLL. A VAF-only fit under the same area convention is a comparator, not a replacement for joint inference. Profile alpha while optimizing theta and vice versa; report flat or boundary-limited support rather than unsupported finite intervals. Numerical bounds and fixed-parameter ranges are not confidence intervals.

\section{Exact stochastic simulation experiment}
The prescribed fractions refer to mutually exclusive $C$ and $D$ counts at day 0:
\begin{center}\begin{tabular}{lrrr}\toprule Initial $D/N$ & $W(0)$ & $C(0)$ & $D(0)$\\\midrule
0.01 & 97000 & 2000 & 1000\\
0 & 98000 & 2000 & 0\\\bottomrule\end{tabular}\end{center}
The corresponding founder VAFs are 1.5\% and 1\%, not 2\%. Simulate to day 1460, saving daily integer counts plus exact first-passage times and replicate seeds. Use
\[
\alpha\in\{2,4,8\},\qquad\theta\in\{2,7,21\}\ \mathrm{days}.
\]
Nine gamma pairs, two initial-D settings and 200 replicates yield 3600 primary trajectories. Also simulate 200 replicates per initial-D setting at the fitted reference and each compatible fitted sensitivity parameter set. The primary nine-point grid is not crossed again with every sensitivity assumption.

The native C++ simulator uses exact time-inhomogeneous event thinning. It analytically skips type-preserving Moran replacements, uses a state-specific bound valid for all subsequent fitness times, and updates W/C/D jointly. For fixed state $n$, use upper bounds on newborn numerators and the lower bound $F\ge N+s_DD$ to bound the sum of off-diagonal rates. Rejected proposals do not change the state. This changes computational efficiency, not the Moran law. Continue paths after threshold crossing.

The four primary figures are single-mutant C (blue) and double-hit D (red), separately for each initial-D setting. Each 3-by-3 grid panel contains 200 curves with opacity 0.035 and a bold arithmetic mean over all replicates, including zeros/extinctions. Show full four-year time and an additional early-time view. Preserve raw counts, compressed daily arrays, seeds, first-passage times, null-proposal counts and runtime metadata.

\section{Risk endpoints, interpretation and validation}
Define $T_*=\inf\{t\ge0:D(t)\ge K\}$. Evaluate cumulative probabilities at days 365, 730, 1095 and 1460, conditional on initial state and fixed parameters. When clinical baseline precedes infusion, condition on an initially event-free day-0 state. The strict 1\% rule gives $K=1001$; the simulation starting at $D=1000$ is one cell below threshold and describes existing-double-hit dynamics, not healthy-baseline incidence.

Use forward first-passage probabilities for rare risks, with simulations as checks. With 200 replicates the empirical event-proportion increment is 0.005; zero events do not establish zero risk. Separate numerical error, parameter uncertainty and fixed-parameter scenario variation. Do not equate a 1\% double-hit threshold with a validated clinical diagnosis or generalize an enriched cohort's conditional fit to all CAR-T recipients.

Validation includes: fixed-area quadrature across the gamma grid; pre-CAR neutrality; no second-hit events when $\mu=D_{\mathrm{init}}=0$; exact population conservation and seed uniqueness; neutral Monte Carlo moments; independent tiny-event matrix-exponential calculations; small-N full Moran versus the inference approximation; numerical cap/tolerance refinements; and observed-data residuals. A successful numerical check does not resolve clinical ascertainment, founder identity, measurement mapping or competing death.

\section{Execution and deliverables}
The code is authored first, followed by this English TeX protocol, then numerical execution. Run validation, the reference joint fit, all compatible sensitivity refits, exact Monte Carlo grids/fitted-point trajectories, first-passage risks and profile likelihoods. Checkpoints permit resumption without repeating finished work. Export the English report PDF, four primary blue/red trajectory figures, fitted-point comparisons, result workbook, CSVs, raw compressed counts and reproducible code. Record configuration, inputs, software versions, timings and seeds. Preserve previous studies separately.

\IfFileExists{executed_results.tex}{\clearpage\section{Executed results}
\textbf{Analysis scope.} The source-record joint cohort contains 14 patients, 42 follow-up VAF observations, 2 exact tMN times (days 127 and 589), 1 source-supported day-1460 censor record and 11 unknown clinical outcomes. The latter contribute VAF likelihood only. These are exploratory, source-cohort conditional estimates, not calibrated population incidence.

\subsection{Reference joint fit}
The reference fit gives $\alpha=1.01$ and $\theta=0.25$ days; the peak occurs at day $0.0025$ with additional fitness $3.8035$. The integrated additional advantage is fixed at 1 day. The joint negative log likelihood is $510.564$, decomposed into VAF $291.916$ and clinical-given-VAF $218.648$. The cap-refinement difference is $3.02e-09$ log units; the tolerance-refinement difference is $2.41e-07$. The optimum reaches a numerical search boundary. Optimizer status: True.
Both parameters reach their lower search limits. The resulting peak precedes the earliest post-CAR VAF observation (day 1); its precise sub-day timing is not empirically resolved and should not be interpreted as a measured biological peak.
\begin{center}\includegraphics[width=\textwidth]{results/figures/joint_fit_diagnostics.pdf}\end{center}
\subsection{Clinical likelihood contributions}
\begin{center}\begin{tabular}{lrrr}\toprule
Patient & Outcome & VAF NLL & Clinical-given-VAF NLL\\\midrule
Patient\_25 & censor & $-3.632\times10^{-17}$ & $1.414\times10^{-8}$ \\
Patient\_103 & exact & $1.29\times10^{-16}$ & $170.35$ \\
Patient\_197 & exact & $11.076$ & $48.303$ \\
\bottomrule\end{tabular}\end{center}
\begin{samepage}\subsection{Standardized four-year probabilities}
Both scenarios start with $C/N=0.02$; $D/N$ is stated separately. Numbers are probabilities, not percentages. The $D/N=0.01$ scenario starts one cell below the strict 1\% threshold and therefore describes existing double-hit dynamics.
\begin{center}\begin{tabular}{lrrrr}\toprule
Initial $D/N$ & Year 1 & Year 2 & Year 3 & Year 4\\\midrule
0 & $1.011\times10^{-32}$ & $5.468\times10^{-17}$ & $1.222\times10^{-11}$ & $5.652\times10^{-9}$ \\
0.01 & $1-1.319\times10^{-9}$ & $1-2.197\times10^{-10}$ & $1-5.224\times10^{-11}$ & $1-1.493\times10^{-11}$ \\
\bottomrule\end{tabular}\end{center}\end{samepage}
\subsection{Fixed-parameter sensitivity}
Every compatible scenario re-estimates the same shared alpha/theta pair using the full joint likelihood. The no-mutation/no-initial-D scenario is structurally incompatible with observed tMN and is recorded as such, without a probability floor. Varying $b$ holds $s_D$ fixed. Full numeric results, including likelihood components and numerical checks, are in the workbook and CSV files.
Changing the observation SD changes the observation model and its likelihood normalization; do not rank biological assumptions solely by cross-SD NLL. The fixed gamma area is over $[0,\infty)$; the area actually falling inside four years is recorded separately.
\begin{longtable}{lrrrr}\toprule
Scenario & $\alpha$ & $\theta$ (days) & Joint NLL & Boundary\\\midrule\endhead
reference & 1.01 & 0.25 & 510.56 & True \\
D\_init\_1 & 1.01 & 0.25 & 492.69 & True \\
D\_init\_10 & 1.01 & 0.25 & 466.65 & True \\
D\_init\_100 & 1.01 & 0.25 & 389.61 & True \\
area\_days\_0.3 & 1.01 & 0.25 & 514.17 & True \\
area\_days\_3 & 1.541 & 0.25 & 501.41 & True \\
b\_0.00357143 & 50 & 1.87 & 1725.1 & True \\
b\_0.142857 & 1.01 & 0.25 & 349.43 & True \\
mu\_0 & \multicolumn{4}{l}{Structurally incompatible} \\
mu\_0.0001 & 1.01 & 0.25 & 500.28 & True \\
mu\_1e-05 & 1.01 & 0.25 & 505.86 & True \\
mu\_1e-07 & 1.01 & 0.25 & 515.18 & True \\
old\_rates\_fixed\_area & 1.01 & 0.25 & 348.76 & True \\
s\_D\_0 & 1.01 & 0.25 & 530.69 & True \\
s\_D\_0.01 & 1.01 & 0.25 & 520.36 & True \\
s\_D\_0.05 & 1.01 & 0.25 & 484.28 & True \\
sigma\_vaf\_0.1 & 1.01 & 0.25 & 2107.4 & True \\
sigma\_vaf\_0.5 & 50 & 1.98 & 303.92 & True \\
threshold\_fraction\_0.005 & 1.01 & 0.25 & 413.09 & True \\
threshold\_fraction\_0.02 & 1.01 & 0.25 & 706.59 & True \\
vaf\_only\_comparator & 1.01 & 0.25 & 291.92 & True \\
\bottomrule\end{longtable}
\clearpage\subsection{Uncertainty and validation}
Profile likelihood results optimize the other parameter at each fixed value. A flat or boundary-limited profile means that a finite two-sided interval is not established by these data. Fixed-parameter scenario variation is not a confidence interval. Initial validation includes fixed-area integrals, an independently exponentiated tiny-event generator, a small-N full Moran comparison and neutral Monte Carlo moment checks. Numerical agreement does not establish the clinical meaning of the threshold or representative sampling.
\begin{center}\begin{tabular}{lrr}\toprule Profile & Fixed value & $\Delta$ NLL\\\midrule
alpha\_profile & 1.01 & 0 \\
alpha\_profile & 2 & 0.3282 \\
alpha\_profile & 10 & 0.968 \\
alpha\_profile & 50 & 1.046 \\
theta\_profile & 0.25 & 0 \\
theta\_profile & 2 & 1.319 \\
theta\_profile & 21 & 1.766 \\
theta\_profile & 1460 & 4.899 \\
\bottomrule\end{tabular}\end{center}
\paragraph{Monte Carlo cross-check.} Exact Moran endpoint means and first-passage proportions are compared with the inference approximation at three grid points and both initial-D settings. Rare double-hit means with $D(0)=0$ have heavy sampling tails; an empirical SE from 200 paths can substantially understate uncertainty when large clones are not sampled. The cross-check therefore also supplies the D-mean Monte Carlo SE predicted by the forward distribution. This is distinct from clinical parameter uncertainty. See \texttt{simulation\_crosscheck.csv}.
\clearpage\section{Exact stochastic Moran trajectories}
\subsection*{Single mutant C: blue; initial $D/N=0.01$}
\includegraphics[width=\textwidth]{results/figures/single_mutant_initial_D_0.01.pdf}
Each panel contains 200 complete random trajectories and their arithmetic mean. Area, not peak height, is held fixed. Trajectories continue after threshold crossing.\clearpage
\subsection*{Double hit D: red; initial $D/N=0.01$}
\includegraphics[width=\textwidth]{results/figures/double_hit_initial_D_0.01.pdf}
Each panel contains 200 complete random trajectories and their arithmetic mean. Area, not peak height, is held fixed. Trajectories continue after threshold crossing.\clearpage
\subsection*{Single mutant C: blue; initial $D/N=0$}
\includegraphics[width=\textwidth]{results/figures/single_mutant_initial_D_0.pdf}
Each panel contains 200 complete random trajectories and their arithmetic mean. Area, not peak height, is held fixed. Trajectories continue after threshold crossing.\clearpage
\subsection*{Double hit D: red; initial $D/N=0$}
\includegraphics[width=\textwidth]{results/figures/double_hit_initial_D_0.pdf}
Each panel contains 200 complete random trajectories and their arithmetic mean. Area, not peak height, is held fixed. Trajectories continue after threshold crossing.\clearpage
\includegraphics[width=\textwidth]{results/figures/joint_optimum_trajectories.pdf}
}{\textit{The protocol is prepared before numerical execution; the results section will be appended after the confirmed pipeline finishes.}}
\end{document}
